% --------------------------------------------------
% Glossar-Einträge
% --------------------------------------------------
\newglossaryentry{autowaiting}{
    name={Auto-Waiting},
    description={Ein Mechanismus, bei dem das Test-Framework automatisch auf den Zustand von DOM-Elementen wartet, bevor eine Aktion ausgeführt wird}
}

\newglossaryentry{boilerplate}{
    name={Boilerplate-Code},
    description={Wiederkehrende Code-Fragmente, die mit wenig oder gar keiner Änderung an mehreren Stellen im Testcode wiederverwendet werden müssen}
}

\newglossaryentry{chaining}{
    name={Chaining},
    description={Eine Programmiertechnik, bei der mehrere Methodenaufrufe direkt aneinandergereiht werden, was zu einer kompakteren Syntax führt}
}

\newglossaryentry{e2etest}{
    name={End-to-End-Test},
    description={Ein Testverfahren, das das gesamte Softwaresystem über alle Schichten hinweg unter realitätsnahen Bedingungen aus Sicht des Endanwenders prüft}
}

\newglossaryentry{e2etests}{
    name={End-to-End-Tests},
    description={Mehrere Testverfahren, die das gesamte Softwaresystem über alle Schichten hinweg unter realitätsnahen Bedingungen aus Sicht des Endanwenders prüft}
}

\newglossaryentry{flakytest}{
    name={Flaky Test},
    description={Ein unzuverlässiger automatisierter Test, der bei gleichem Systemzustand inkonsistente Ergebnisse (Erfolg oder Fehlschlag) liefert}
}

\newglossaryentry{regkrit}{
    name={regressionskritischen Testszenarien},
    description={Anzahl von Tests zur Überprüfung, ob bestehende Funktionalitäten der Software nach Änderungen weiterhin fehlerfrei arbeiten}
}

\newglossaryentry{timetravel}{
    name={Time-Travel-Funktion},
    description={Ein Feature (z B in Cypress), das den exakten visuellen Zustand der Weboberfläche zu jedem vergangenen Testschritt speichert und nachträglich abrufbar macht}
}

\newglossaryentry{webdriver}{
    name={WebDriver},
    description={Ein Standard-Protokoll zur nativen Fernsteuerung von Webbrowsern, das primär von Selenium genutzt wird}
}

\newglossaryentry{rueckwaertskompatibilitaet}{
    name={Rückwärtskompatibilität},
    description={Die Eigenschaft einer Software oder eines Frameworks, mit älteren Versionen, Schnittstellen oder Laufzeitumgebungen (z\,B älteren Java-Versionen) ohne funktionale Einschränkungen fehlerfrei zusammenzuarbeiten}
}